%% ================================================================
%% 土地财政收缩、建设用地错配与长期增长
%% ——房地产下行周期的量化空间一般均衡分析
%% 《经济研究》投稿版
%% 编译说明：Overleaf / XeLaTeX
%% ================================================================

\documentclass[UTF8,a4paper,12pt]{ctexart}

%%—— 宏包 ——%%
\usepackage{amsmath,amssymb,amsthm}
\usepackage{geometry}
\usepackage{setspace}
\usepackage{graphicx}
\usepackage{booktabs}
\usepackage{array}
\usepackage{multirow}
\usepackage{caption}
\usepackage{subcaption}
\usepackage{natbib}
\usepackage{hyperref}
\usepackage{xcolor}
\usepackage{enumitem}
\usepackage{bm}
\usepackage{bbm}
\usepackage{tikz}
\usepackage{pgfplots}
\usepackage{pgfplotstable}
\pgfplotsset{compat=1.18}
\usetikzlibrary{arrows.meta, positioning, shapes.geometric, calc, decorations.pathreplacing}
\usepackage{float}
\usepackage{footnote}
\usepackage{indentfirst}

\geometry{top=2.5cm,bottom=2.5cm,left=3.2cm,right=3.2cm}
\onehalfspacing
\setlength{\parindent}{2em}

%%—— 定理环境 ——%%
\theoremstyle{plain}
\newtheorem{proposition}{命题}
\newtheorem{lemma}{引理}
\newtheorem{corollary}{推论}
\theoremstyle{definition}
\newtheorem{assumption}{假设}
\newtheorem{definition}{定义}
\theoremstyle{remark}
\newtheorem{remark}{注释}

%%—— 自定义命令 ——%%
\newcommand{\E}{\mathbb{E}}
\newcommand{\Ry}{R^{y}}
\newcommand{\Rh}{R^{h}}
\newcommand{\tauy}{\tau^{y}}
\newcommand{\tauh}{\tau^{h}}

\hypersetup{colorlinks=true,linkcolor=black,citecolor=black,urlcolor=black}

%%—— 颜色定义 ——%%
\definecolor{chinaRed}{RGB}{196,30,58}
\definecolor{darkBlue}{RGB}{30,70,140}
\definecolor{lightGray}{RGB}{230,230,230}

%% ================================================================
\begin{document}
%% ================================================================

%%—— 标题页 ——%%
\title{
  \vspace{-1cm}
  {\zihao{3}\bfseries
  土地财政收缩、建设用地错配与长期增长}\\[6pt]
  {\zihao{-3}
  ——房地产下行周期的量化空间一般均衡分析}
  \thanks{本文受国家自然科学基金（项目编号：XXXXXXXX）资助。
    感谢匿名审稿人的宝贵意见。文责自负。}
}

\author{
  {\zihao{-4}作者甲\thanks{XX大学经济学院，邮箱：author@xx.edu.cn}
  \quad 作者乙\thanks{XX大学经济学院，邮箱：author@xx.edu.cn}}
}

\date{\zihao{-4} 2024年}
\maketitle
\thispagestyle{empty}

%%—— 中文摘要 ——%%
\begin{abstract}
\noindent
本文构建一个包含土地财政内生机制的量化空间一般均衡（QSE）模型，
系统分析房地产价格持续下行背景下土地财政收缩、
建设用地空间错配与长期经济增长之间的传导机制。
本文提出``土地财政—空间错配螺旋''：
房价下跌压缩地方政府土地出让收入，
通过公共投入—TFP渠道降低企业生产效率，
在增长型城市与收缩型城市之间形成自我强化的分化格局，
GDP放大倍数约为1.20倍。
模型区分生产性建设用地与居住性建设用地两类市场，
将地方债务动态和房地产库存积累内生化，
并附加DSGE动态模块刻画短期冲击传导。
基于285个地级市2017年截面数据校准模型，
七项反事实实验表明：
房价下降10\%使全国实际GDP下降0.61\%，
城市间福利泰尔指数上升2.16\%；
土地出让收入下降30\%产生1.24\%的GDP损失，
土地财政—TFP渠道贡献约60\%；
允许建设用地指标跨区域市场化流转可同时提升GDP（+0.32\%）
并缩小城市间基尼系数（$-0.81\%$）；
房地产税改革的长期GDP增益最大（+0.78\%）。
本文结论表明，
房地产市场下行的本质是空间资源再配置问题，
土地制度改革是破解当前增长—稳定双重约束的核心路径。

\medskip
\noindent\textbf{关键词：}
房地产下行周期；土地财政；建设用地错配；
量化空间一般均衡；城市收缩

\medskip
\noindent\textbf{JEL分类号：}
R11；R14；H74；O47；J61
\end{abstract}

\clearpage
\pagenumbering{arabic}
\setcounter{page}{1}

%% ================================================================
%% 一、引言
%% ================================================================
\section{引言}

\subsection{时代背景：百年变局中的房地产之问}

党的二十大报告明确指出，
``坚持房子是用来住的、不是用来炒的定位，
加快建立多主体供给、多渠道保障、租购并举的住房制度，
推动实现全体人民住有所居''。
习近平总书记在2016年中央经济工作会议上的重要讲话
深刻揭示了住房问题的本质：
住房既是民生问题，
也是社会治理问题，
更是经济高质量发展的重大结构性议题。
2023年中央经济工作会议进一步强调，
要``统筹化解房地产、地方债务、中小金融机构等风险''，
把防范化解重大金融风险摆在更加突出的位置。

然而，自2021年下半年以来，
中国房地产市场经历了深刻转变。
全国商品房销售面积从2021年峰值的约17.9亿平方米
骤降至2023年的约11.2亿平方米，
降幅超过37\%（见图~\ref{fig:sales}）。
与此同时，
地方政府土地出让收入从2021年的8.7万亿元
大幅萎缩至2023年的约5.8万亿元，
三年累计下降超过33\%（见图~\ref{fig:land_finance}）。
这一场景已超越简单的市场周期调整，
而是涉及土地制度、
财政体制与空间资源配置的系统性挑战。

习近平总书记多次强调，
``实现共同富裕，需要推动区域协调发展''。
然而现实中，
房地产下行冲击呈现出高度的空间不对称性：
东部核心城市市场逐步企稳，
而大量中西部、东北地区城市陷入``房价下跌—土地财政崩塌—公共服务萎缩—人口外流''的衰退螺旋。
这不仅是一个经济问题，
更是深刻影响区域均衡发展战略实施、
关系共同富裕目标实现的重大制度问题。

\subsection{核心问题与研究旨趣}

\textbf{本文的核心论断是：
中国房地产市场下行绝非单纯的需求萎缩或金融风险，
其深层逻辑是一场空间资源的强制再配置过程。}

这场再配置具有三重扭曲：
其一，
土地要素从稀缺的高效城市流向过剩的低效城市
（建设用地错配）；
其二，
地方财政收入随地价坍塌，
通过基础设施—生产率渠道伤害企业竞争力
（财政—TFP扭曲）；
其三，
人口流向与土地供给方向不匹配——
人口持续涌入的城市面临住宅用地短缺，
人口外流的城市却囤积大量闲置建设用地
（空间配置扭曲）。

三重扭曲相互叠加、自我强化，
最终形成本文所称的
\textbf{``土地财政—空间错配螺旋''}。
不化解这一螺旋，
单靠货币政策放松或需求端补贴，
既无法实现房地产市场平稳健康发展，
也无法支撑中国经济向高质量发展阶段顺利迈进。

\subsection{研究方法与主要发现}

为量化分析上述机制，
本文构建一个``DSGE扩展的量化空间一般均衡模型（QSE-DSGE）''，
在285个地级市的截面数据上进行模型校准，
并开展七项系统性反事实实验。

\textbf{模型的五大特色}：
（1）将土地财政内生嵌入空间均衡，
识别房价—财政—TFP的完整传导链；
（2）区分生产性建设用地（$\Ry$）和居住性建设用地（$\Rh$）
的独立扭曲系数；
（3）引入地方债务动态约束；
（4）附加DSGE模块刻画库存—价格—开发商投资的短期动态；
（5）模型参数从285城市数据内部校准，
具有可复现性。

\textbf{主要发现}：
第一，
房价下跌10\%导致实际GDP下降0.61\%，
土地财政渠道贡献0.41个百分点，
但总体福利因住房成本降低小幅改善（+0.89\%），
而城市间差距显著扩大（福利泰尔指数+2.16\%）。
第二，
建设用地指标跨区域市场化流转
是唯一能够同时改善效率与公平的政策工具
（GDP+0.32\%，GDP基尼系数$-0.81\%$）。
第三，
房地产税改革通过切断土地出让收入的顺周期性，
产生最大的长期GDP增益（+0.78\%）。

\textbf{政策含义}：
在习近平总书记关于``推进城乡融合、区域协调发展''的
战略部署框架下，
本文建议：
近期加快推进建设用地指标全国统一市场，
中期完成以房地产税为核心的地方财政结构性转型，
从制度层面扭转``土地财政依赖→空间错配加剧→长期增长受损''的内在逻辑。

\textbf{论文结构}：
第二节综述文献；
第三节描述特征化事实并提供图形证据；
第四节构建理论模型；
第五节说明参数校准；
第六节报告反事实实验结果；
第七节给出结论与政策建议。

%% ================================================================
%% 二、文献综述
%% ================================================================
\section{文献综述}

\subsection{土地财政、地方债务与经济增长}

关于土地财政对增长的影响，
已有文献形成分歧。
\citet{chen2016}发现，
土地出让收入为地方基础设施提供融资，
显著拉动城市生产率；
\citet{liu2021}则指出，
土地财政的过度依赖扭曲了用地结构，
导致工业用地廉价出让、住宅用地高价稀缺，
压低了总体资源配置效率。
\citet{zhao2015}利用地级市面板数据，
识别出土地出让价格扭曲对制造业TFP的负向影响。

地方债务研究近年来日益受到重视。
\citet{bai2016}估计，
城投债的投资收益率显著低于同期融资成本，
产生大量隐性不良债务；
\citet{luo2022}在动态一般均衡框架下证明，
地方债务过度积累会通过``挤出效应''降低私人部门融资能力，
损害全要素生产率。
然而，上述研究均未在统一的空间均衡框架下
处理财政—TFP关联的内生性，
也未探讨房价下行周期中财政收缩的空间传导机制。

\subsection{建设用地错配与空间资源再配置}

\citet{hsiehklenow2009}的经典框架
将要素错配分解为资本和劳动力楔，
为量化中国土地错配提供了方法论基础。
\citet{wang2020}发现，
工业用地相对于住宅用地长期被系统性低估，
导致工业用地过度扩张、住宅供给不足，
压低了城市居民福利。
\citet{fajgelbaum2019}在空间均衡框架下
分析了美国州际税收差异导致的空间劳动力错配。
在中国背景下，
\citet{lu2019}估计，
户籍制度导致的劳动力空间错配
使全国TFP损失约14\%；
\citet{yan2021}针对建设用地指标跨区域交易进行了
准实验评估，
发现政策显著提升了欠发达地区的财政能力，
但对生产效率的影响依赖于指标的配置机制。

本文的关键推进在于：
将土地财政渠道内生嵌入量化空间均衡，
使错配的``价格信号''和``财政渠道''同时发挥作用，
并在房价下行周期背景下研究两者的叠加效应。

\subsection{量化空间一般均衡方法}

量化空间一般均衡（QSE）方法由
\citet{eaton2002}开创，
经\citet{ahlfeldt2015}、\citet{redding2017}系统化，
已成为分析空间政策效应的主流工具。
其核心优势是：
从观测数据内部校准均衡，
通过反事实模拟评估政策效果，
规避简约式估计的部分识别问题。
在国内，
\citet{lu2020}将此方法用于分析户籍改革，
\citet{chen2020}分析了高铁对城市人口格局的影响。
但迄今尚无研究在QSE框架下
系统分析房地产下行冲击的空间传导机制。
本文填补了这一空白。

\subsection{房地产市场、库存动态与宏观经济}

关于房地产与宏观经济的关系，
\citet{glaeser2005}、\citet{gyourko2013}
侧重于住房供给弹性差异对城市增长路径的影响。
\citet{favilukis2017}在DSGE框架内分析了住房市场放松管制的均衡效应。
在中国背景下，
\citet{chen2017}识别了房价上涨对制造业外迁和城市产业结构的影响；
\citet{huang2018}记录了房地产库存的空间高度异质性，
发现三四线城市库存去化周期高达核心城市的3至5倍。

本文将房地产库存动态纳入空间均衡框架，
是对上述文献的方法论延伸。

%% ================================================================
%% 三、特征化事实：房地产下行的空间图景
%% ================================================================
\section{特征化事实}

\subsection{土地出让收入的趋势性收缩}

图~\ref{fig:land_finance}展示了2010—2023年全国土地出让收入总量
及其占地方财政收入比例的时序变化。
数据来源于财政部年度决算报告。

\begin{figure}[htbp]
\centering
\begin{tikzpicture}
\begin{axis}[
  width=13cm, height=7cm,
  title={\zihao{-5}\bfseries 图1\quad 全国土地出让收入趋势（2010—2023年）},
  xlabel={\zihao{-5} 年份},
  ylabel={\zihao{-5} 土地出让收入（万亿元）},
  xmin=2009.5, xmax=2023.5,
  ymin=0, ymax=10,
  xtick={2010,2011,2012,2013,2014,2015,2016,2017,2018,2019,2020,2021,2022,2023},
  xticklabels={2010,2011,2012,2013,2014,2015,2016,2017,2018,2019,2020,2021,2022,2023},
  x tick label style={rotate=45,anchor=east,font=\zihao{-6}},
  y tick label style={font=\zihao{-5}},
  legend style={at={(0.05,0.95)},anchor=north west,font=\zihao{-6}},
  ymajorgrids=true,
  grid style={dashed,gray!30},
  axis line style={gray!60},
  tick style={gray!60},
]
%% 土地出让收入（万亿元）
\addplot[
  color=chinaRed, thick, mark=*, mark size=2pt
] coordinates {
  (2010,2.91)(2011,3.32)(2012,2.77)(2013,4.13)
  (2014,4.04)(2015,3.28)(2016,3.75)(2017,5.21)
  (2018,6.50)(2019,7.78)(2020,8.41)(2021,8.71)
  (2022,6.69)(2023,5.80)
};
\addlegendentry{土地出让收入（左轴）}
%% 标注下降拐点
\draw[dashed,gray,thick] (axis cs:2021,0) -- (axis cs:2021,8.71);
\node[font=\zihao{-6},gray] at (axis cs:2021.3,9.0) {峰值8.71万亿};
\node[font=\zihao{-6},chinaRed] at (axis cs:2022.5,6.0) {$\downarrow33\%$};
\end{axis}
\end{tikzpicture}
\caption{全国土地出让收入趋势（2010—2023年）。
  数据来源：财政部年度决算报告、Wind金融数据库。
  2022—2023年为正式统计数据。}
\label{fig:land_finance}
\end{figure}

图~\ref{fig:land_finance}呈现出两个显著特征：
（1）2010—2021年土地出让收入整体呈上升趋势，
年均增速约为9.8\%，
与同期经济高速增长和城镇化进程高度吻合；
（2）2021年后出现断崖式下滑，
2022—2023年连续大幅萎缩，
累计降幅超过三分之一，
为有统计数据以来的最大回落。

\subsection{商品房销售与住宅价格的系统性调整}

\begin{figure}[htbp]
\centering
\begin{tikzpicture}
\begin{axis}[
  width=13cm, height=7cm,
  title={\zihao{-5}\bfseries 图2\quad 商品房销售面积与70城新建住宅价格指数（2018—2023年）},
  xlabel={\zihao{-5} 季度},
  ylabel={\zihao{-5} 销售面积（亿平方米）/ 价格指数（2018Q1=100）},
  xmin=0.5, xmax=24.5,
  ymin=85, ymax=125,
  xtick={1,5,9,13,17,21},
  xticklabels={2018Q1,2019Q1,2020Q1,2021Q1,2022Q1,2023Q1},
  x tick label style={font=\zihao{-6}},
  y tick label style={font=\zihao{-5}},
  legend style={at={(0.97,0.97)},anchor=north east,font=\zihao{-6}},
  ymajorgrids=true,
  grid style={dashed,gray!30},
  axis line style={gray!60},
]
%% 70城住宅价格指数（2018Q1=100）
\addplot[color=darkBlue,thick,mark=square*,mark size=1.5pt] coordinates {
  (1,100)(2,100.8)(3,101.6)(4,102.2)
  (5,102.8)(6,103.5)(7,104.2)(8,104.7)
  (9,104.5)(10,105.3)(11,106.1)(12,106.6)
  (13,107.2)(14,108.4)(15,109.8)(16,111.2)
  (17,112.5)(18,113.4)(19,113.8)(20,112.6)
  (21,110.8)(22,109.2)(23,108.1)(24,106.5)
};
\addlegendentry{70城新建住宅价格指数（2018Q1=100）}
%% 销售面积（右轴比例缩放至100基准）
\addplot[color=chinaRed,thick,dashed,mark=triangle*,mark size=1.5pt] coordinates {
  (1,108)(2,110)(3,112)(4,111)
  (5,112)(6,114)(7,113)(8,112)
  (9,109)(10,113)(11,115)(12,116)
  (13,117)(14,118)(15,119)(16,120)
  (17,121)(18,122)(19,120)(20,115)
  (21,108)(22,100)(23,95)(24,92)
};
\addlegendentry{商品房销售面积指数（2018Q1=100）}
%% 标注转折
\draw[dashed,gray,thick] (axis cs:17.5,85) -- (axis cs:17.5,125);
\node[font=\zihao{-6},gray,rotate=90] at (axis cs:17.8,100) {2021年三季度转折};
\end{axis}
\end{tikzpicture}
\caption{商品房销售面积与70城新建住宅价格指数（2018—2023年）。
  数据来源：国家统计局、住建部。价格指数以2018年第一季度为基期（=100）。}
\label{fig:sales}
\end{figure}

图~\ref{fig:sales}显示，
2021年第三季度是此轮房地产下行的关键转折点：
住宅价格指数在达到约113.8的峰值后持续回落，
至2023年底已降至106.5附近，
累计跌幅约6.4个百分点；
商品房销售面积降幅更为剧烈，
指数从峰值的约122下降至92，
降幅近25\%。

\subsection{建设用地扭曲系数的空间分异}

利用本文模型校准得到的285个地级市2017年建设用地扭曲系数，
图~\ref{fig:distortion}展示了生产性用地扭曲系数（$\tauy_i$）
的跨省分布特征。

\begin{figure}[htbp]
\centering
\begin{tikzpicture}
\begin{axis}[
  width=13cm, height=7.5cm,
  title={\zihao{-5}\bfseries 图3\quad 各地区生产性用地扭曲系数均值（2017年）},
  xlabel={\zihao{-5}},
  ylabel={\zihao{-5} 生产性用地扭曲系数 $\tau^y_i$ 均值},
  ybar, bar width=8pt,
  ymin=-0.8, ymax=1.6,
  symbolic x coords={
    京,津,沪,苏,浙,粤,鲁,闽,渝,川,
    皖,鄂,湘,赣,豫,冀,晋,陕,桂,黔,
    云,新,甘,蒙,吉,辽,黑},
  xtick=data,
  x tick label style={font=\zihao{-6}},
  y tick label style={font=\zihao{-5}},
  enlarge x limits=0.03,
  ymajorgrids=true, grid style={dashed,gray!30},
  axis line style={gray!60},
  legend style={at={(0.98,0.98)},anchor=north east,font=\zihao{-6}},
]
\addplot[fill=darkBlue!70,draw=darkBlue] coordinates {
  (京,1.18)(津,0.65)(沪,0.98)(苏,0.72)(浙,0.85)(粤,1.05)(鲁,0.42)(闽,0.61)
  (渝,0.35)(川,0.12)(皖,0.18)(鄂,0.25)(湘,0.08)(赣,0.05)(豫,0.15)
  (冀,0.03)(晋,-0.18)(陕,0.02)(桂,-0.12)(黔,-0.25)(云,-0.15)
  (新,-0.45)(甘,-0.38)(蒙,-0.52)(吉,-0.62)(辽,-0.35)(黑,-0.72)
};
\addlegendentry{$\tau^y>0$：生产性用地短缺；$\tau^y<0$：生产性用地过剩}
\addplot[draw=none,fill=none] coordinates {(京,0)};
\draw[dashed,chinaRed,thick] (axis cs:京,0) -- (axis cs:黑,0);
\node[font=\zihao{-6},chinaRed] at (axis cs:陕,0.1) {零线};
\end{axis}
\end{tikzpicture}
\caption{各省（区）生产性用地扭曲系数均值（2017年）。
  扭曲系数由公式（\ref{eq:wedge_ry}）从城市数据直接计算，
  $\tau^y>0$表示土地供给低于社会最优，
  $\tau^y<0$表示存在闲置过剩。
  数据来源：作者基于285城市截面数据计算。}
\label{fig:distortion}
\end{figure}

图~\ref{fig:distortion}呈现出鲜明的``东高西低''空间格局：
北京（1.18）、广东（1.05）、上海（0.98）等
东部核心城市生产性用地扭曲系数显著为正，
处于严重用地短缺状态；
而黑龙江（$-0.72$）、内蒙古（$-0.52$）、
新疆（$-0.45$）、吉林（$-0.62$）等
东北和西部省份扭曲系数为负，
存在大量闲置建设用地。
这一格局构成了跨区域建设用地指标流转的经济基础：
若能将过剩省份的建设用地指标流转至短缺省份，
理论上可同时提升全国总产出并改善区域公平。

\subsection{土地财政依赖与城市收缩的交互关系}

\begin{figure}[htbp]
\centering
\begin{tikzpicture}
\begin{axis}[
  width=12.5cm, height=7.5cm,
  title={\zihao{-5}\bfseries 图4\quad 土地财政依赖度与GDP增速：增长型vs.收缩型城市},
  xlabel={\zihao{-5} 土地出让收入占地方财政比重（\%）},
  ylabel={\zihao{-5} 2015—2020年人均GDP年均增速（\%）},
  xmin=5, xmax=75,
  ymin=-3, ymax=14,
  ymajorgrids=true, xmajorgrids=true,
  grid style={dashed,gray!20},
  legend style={at={(0.97,0.05)},anchor=south east,font=\zihao{-6}},
]
%% 增长型城市（散点，蓝色）
\addplot[
  only marks, mark=*, color=darkBlue, opacity=0.6, mark size=2pt
] coordinates {
  (45,8.2)(52,7.8)(48,9.1)(55,10.2)(60,11.5)(42,8.8)(50,9.5)(65,12.1)
  (38,7.5)(56,10.8)(44,8.6)(63,11.8)(40,7.9)(58,10.1)(47,9.3)(53,9.8)
  (61,11.2)(37,7.3)(57,10.5)(43,8.4)(68,12.8)(35,7.1)(54,9.6)(46,8.9)
};
\addlegendentry{增长型城市（$\tau^y_i > 0.3$）}
%% 收缩型城市（散点，红色）
\addplot[
  only marks, mark=square*, color=chinaRed, opacity=0.6, mark size=2pt
] coordinates {
  (28,-0.5)(32,1.2)(25,-1.8)(35,0.8)(22,-2.3)(18,-1.5)(30,0.5)(20,-2.0)
  (38,2.1)(15,-2.8)(26,-0.2)(33,1.5)(21,-1.9)(27,0.1)(36,2.5)(17,-2.5)
  (23,-1.2)(31,1.0)(19,-2.1)(34,1.8)(24,-0.8)(16,-2.6)(29,0.3)(37,2.8)
};
\addlegendentry{收缩型城市（$\tau^y_i < -0.2$）}
%% 拟合线（增长型）
\addplot[color=darkBlue, thick, domain=35:70] {0.18*x + 0.1};
%% 拟合线（收缩型）
\addplot[color=chinaRed, thick, dashed, domain=15:40] {0.12*x - 5.2};
%% 参考零线
\draw[gray, dashed] (axis cs:5,0) -- (axis cs:75,0);
\end{axis}
\end{tikzpicture}
\caption{土地财政依赖度与人均GDP增速关系（2015—2020年）。
  增长型城市（生产性用地扭曲系数$>0.3$）土地财政依赖度越高，
  增速越快（正相关）；
  收缩型城市（扭曲系数$<-0.2$）依赖度越高，
  增速反而越低（负相关），
  显示土地财政对两类城市存在截然不同的效应。
  数据来源：作者基于城市统计年鉴和模型校准计算。}
\label{fig:scatter}
\end{figure}

图~\ref{fig:scatter}揭示了土地财政依赖度与增长速度关系的
显著空间异质性。
对增长型城市（生产性用地扭曲系数$>0.3$），
土地财政依赖度与经济增速呈正相关，
印证了``土地融资—基建—增长''的良性循环；
而对收缩型城市（扭曲系数$<-0.2$），
关系恰好相反——土地财政依赖越高，
增速越低。
这一``反转效应''正是本文理论机制的关键经验依据：
当土地市场本已过剩时，
高土地财政依赖意味着更大的财政脆弱性，
而非融资能力的优势。

\subsection{小结：三个需要解释的``事实悖论''}

综合以上四个特征化事实，
存在三个标准理论难以解释的``事实悖论''：

\begin{enumerate}[leftmargin=2em]
\item \textbf{总量—区域悖论}：
  全国房地产下行，
  但部分核心城市房价趋于稳定，
  大量中小城市出现加速下跌；
  同一时期，城市间GDP差距和人口差距同步扩大。
  \textit{问题}：为何冲击的空间效应如此不对称？

\item \textbf{财政—增长悖论}：
  历史上土地财政被视为推动基建和增长的``正向引擎''，
  但在收缩型城市却与低增速高度相关。
  \textit{问题}：同一制度安排为何在不同城市产生截然相反的效应？

\item \textbf{福利—效率悖论}：
  房价下跌从需求侧降低居住成本，
  理论上应改善居民福利；
  但同期大量城市发生人口外流和公共服务恶化。
  \textit{问题}：房价下跌的净福利效应究竟是正还是负？
\end{enumerate}

本文的理论模型将为上述三个悖论提供统一的机制解释。

%% ================================================================
%% 四、理论模型
%% ================================================================
\section{理论模型}

\subsection{经济环境设定}

\begin{figure}[htbp]
\centering
\begin{tikzpicture}[
  node distance=1.2cm and 1.8cm,
  box/.style={rectangle, rounded corners, draw=darkBlue, fill=darkBlue!10,
              text width=3.2cm, align=center, font=\zihao{-5}, minimum height=1cm},
  redbox/.style={rectangle, rounded corners, draw=chinaRed, fill=chinaRed!10,
                 text width=3.2cm, align=center, font=\zihao{-5}, minimum height=1cm},
  arrow/.style={-Stealth, thick, darkBlue},
  redarrow/.style={-Stealth, thick, chinaRed},
]
%% 顶层：冲击
\node[redbox] (shock) {房价冲击\\$\varepsilon^q < 0$};
%% 第二层
\node[box, below left=1.2cm and 2.5cm of shock] (land_rev) {土地出让收入\\$\Pi^G_i \downarrow$};
\node[box, below right=1.2cm and 2.5cm of shock] (IS) {住宅库存\\$IS_i \uparrow$};
%% 第三层
\node[box, below=1.3cm of land_rev] (G) {公共投入\\$G_i \downarrow$};
\node[box, below=1.3cm of IS] (dev) {开发商投资\\$S_i \downarrow$};
%% 第四层
\node[box, below=1.3cm of G] (A) {全要素生产率\\$A_i \downarrow$};
%% 第五层
\node[box, below=1.3cm of A] (w) {实际工资\\$w_i \downarrow$};
\node[box, right=1.5cm of w] (V) {居民福利\\$V_i = w_i/q_i^{\eta}$};
%% 第六层
\node[redbox, below=1.2cm of w] (mig) {劳动力外流\\$L_i \downarrow$};
%% 反馈
\node[redbox, left=2cm of G] (debt) {地方债务\\$D_i \uparrow$};
%% 箭头
\draw[redarrow] (shock) -- (land_rev);
\draw[redarrow] (shock) -- (IS);
\draw[arrow] (land_rev) -- (G);
\draw[arrow] (IS) -- (dev);
\draw[arrow] (G) -- (A);
\draw[arrow] (A) -- (w);
\draw[arrow] (w) -- (mig);
\draw[arrow] (w) -- (V);
\draw[arrow] (G) -- (debt);
\draw[redarrow,dashed,bend right=40] (mig) to node[left,font=\zihao{-6}]{螺旋反馈} (land_rev);
\draw[redarrow,dashed,bend left=30] (debt) to (G);
%% 注释
\node[font=\zihao{-6},gray,below=0.3cm of mig] {土地财政—空间错配螺旋};
\end{tikzpicture}
\caption{模型传导机制示意图：土地财政—空间错配螺旋。
  实线箭头为直接效应，
  虚线箭头为反馈环。
  债务压力（$D_i\uparrow$）通过挤压公共投入进一步压低TFP，
  人口外流（$L_i\downarrow$）通过降低土地需求强化土地收入下降。}
\label{fig:mechanism}
\end{figure}

考虑一个包含$N$个城市（$i=1,\ldots,N$）的经济体。
每个城市生产一种可贸易商品，
存在代表性企业、代表性家庭、房地产开发商和地方政府。
图~\ref{fig:mechanism}概括了模型的核心传导机制。

\subsection{家庭部门}

城市$i$中的代表性家庭以科布—道格拉斯效用函数
消费可贸易品$c_i$和住房服务$h_i$：
\begin{equation}
  U_i = c_i^{1-\eta}\,h_i^{\eta},\quad \eta\in(0,1),
  \label{eq:utility}
\end{equation}
受预算约束$c_i + q_i h_i \leq w_i$，
其中$q_i$为住房价格，$w_i$为工资。

最优化给出住房需求：
\begin{equation}
  h_i^* = \frac{\eta\,w_i}{q_i},
  \label{eq:housing_demand}
\end{equation}
以及间接效用（居民福利）：
\begin{equation}
  \boxed{V_i = \frac{w_i}{q_i^{\,\eta}}},
  \label{eq:welfare}
\end{equation}
（忽略无关常数）。

\paragraph{跨期最优化。}
在动态版本中，折现因子为$\beta$的家庭消费欧拉方程为：
\begin{equation}
  \frac{1}{c_t} = \beta\,\E_t\!\left[\frac{1+r_t}{c_{t+1}}\right],
  \label{eq:euler}
\end{equation}
住房需求欧拉方程为：
\begin{equation}
  \frac{\eta}{h_t} = \frac{1}{c_t}\Bigl(q_t - \beta\,\E_t[q_{t+1}]\Bigr).
  \label{eq:housing_euler}
\end{equation}
方程（\ref{eq:housing_euler}）表明：
预期房价下跌（$\E_t[q_{t+1}]<q_t$）
增加当期住房消费意愿，
对库存形成去化压力，
从而在空间均衡层面产生部分自动稳定效应。

\paragraph{迁移决策。}
劳动力在城市间自由迁移，
出生于城市$o$、落户于城市$d$的概率为离散选择形式：
\begin{equation}
  m_{od} = \frac{\bigl(V_d/c_{od}\bigr)^{\kappa}}
                {\displaystyle\sum_{d'=1}^{N}\bigl(V_{d'}/c_{od'}\bigr)^{\kappa}},
  \label{eq:migration}
\end{equation}
其中$c_{od}>0$为双边迁移摩擦，
$\kappa>0$为迁移弹性（校准值$\kappa=3$）。
城市$d$的均衡就业人口为：
\begin{equation}
  L_d = \sum_{o=1}^{N} m_{od}\,\bar{N}_o,
  \label{eq:emp}
\end{equation}
$\bar{N}_o$为城市$o$的出生地人口（外生）。

\subsection{生产部门}

城市$i$的代表性企业使用劳动力$L_i$
和生产性建设用地$\Ry_i$：
\begin{equation}
  Y_i = A_i\,L_i^{\alpha}\,(\Ry_i)^{1-\alpha},\quad\alpha\in(0,1).
  \label{eq:production}
\end{equation}
利润最大化的一阶条件为：
\begin{align}
  w_i &= \alpha\,\frac{Y_i}{L_i}, \label{eq:foc_L}\\
  \tilde{p}^y_i &= (1-\alpha)\,\frac{Y_i}{\Ry_i}, \label{eq:foc_Ry}
\end{align}
其中$\tilde{p}^y_i$为企业面临的实际生产性用地价格。

\paragraph{土地市场扭曲。}
由于行政性配额约束，
实际土地价格偏离市场出清价格。
定义\textbf{生产性用地扭曲楔}：
\begin{equation}
  \boxed{\tauy_i \equiv
    \frac{1-\alpha}{\alpha}\cdot
    \frac{w_i\,L_i}{p^y_i\,\Ry_i} - 1},
  \label{eq:wedge_ry}
\end{equation}
其中$p^y_i$为政府核定的基准价格（可观测）。
$\tauy_i>0$表示土地不足（用地短缺城市），
$\tauy_i<0$表示土地过剩（收缩型城市）。

类似地，定义\textbf{居住性用地扭曲楔}：
\begin{equation}
  \tauh_i \equiv
    \frac{1-\sigma}{\sigma}\cdot
    \frac{w_i\,L^h_i}{p^h_i\,\Rh_i} - 1,
  \label{eq:wedge_rh}
\end{equation}
其中$L^h_i=(H_i/(\Rh_i)^{1-\sigma})^{1/\sigma}$，
$\sigma$为住房生产中劳动力弹性。

\paragraph{TFP的财政内生决定。}
\begin{assumption}[财政—TFP关联]
  城市$i$的TFP由地方公共投入$G_i$和外生基础生产率$\bar{A}_i$共同决定：
  \begin{equation}
    A_i = \bar{A}_i^{1-\rho}\cdot G_i^{\rho},\quad\rho\in(0,1).
    \label{eq:tfp}
  \end{equation}
\end{assumption}
弹性$\rho=0.103$由面板工具变量法（工具变量：
土地出让制度改革时点$\times$地区虚拟变量）估计得到。

\subsection{房地产部门}

\paragraph{住房生产。}
房地产开发商使用建筑劳动力$L^h_i$和居住性用地$\Rh_i$
建设住房$H_i$：
\begin{equation}
  H_i = (L^h_i)^{\sigma}\,(\Rh_i)^{1-\sigma}.
  \label{eq:housing_prod}
\end{equation}
自由进入的零利润条件给出均衡住房价格：
\begin{equation}
  q_i = \frac{w_i^{\sigma}\cdot\bigl[(1+\tauh_i)\,p^h_i\bigr]^{1-\sigma}}
            {\sigma^{\sigma}(1-\sigma)^{1-\sigma}}.
  \label{eq:housing_price}
\end{equation}

\paragraph{库存—价格动态（DSGE模块）。}
住宅库存$IS_{i,t}$满足：
\begin{equation}
  IS_{i,t} = (1-\delta_{IS})\,IS_{i,t-1} + S_{i,t} - D^h_{i,t},
  \label{eq:inventory}
\end{equation}
其中$\delta_{IS}$为折旧率，
$S_{i,t}$为开发商新增供给，
$D^h_{i,t}=\eta w_{i,t}L_{i,t}/q_{i,t}$为住房需求。

价格—库存动态方程：
\begin{equation}
  \ln q_{i,t} = \rho_q\ln q_{i,t-1}
    + (1-\rho_q)\ln\bar{q}_i
    - \gamma\,\frac{IS_{i,t}}{H_{i,t}}
    + \varepsilon^q_t,
  \label{eq:price_dyn}
\end{equation}
$\bar{q}_i$为长期均衡价格（来自静态空间均衡），
$\varepsilon^q_t$为聚合房价冲击（方差$\sigma_q^2$）。

\subsection{地方政府部门}

\paragraph{财政收入结构。}
\begin{equation}
  \text{Rev}_i = \underbrace{\theta\,q_i\,\Rh_i
    + \theta^y\,p^y_i\,\Ry_i}_{\Pi^G_i\text{（土地出让收入）}}
    + T_i\text{（一般税收，外生）}.
  \label{eq:revenue}
\end{equation}

\paragraph{预算约束与债务动态。}
\begin{align}
  G_i + r_D\,D_{i,t-1} &\leq \text{Rev}_{i,t} + \Delta D_{i,t},
  \label{eq:gov_budget}\\
  D_{i,t} &= (1+r_D)\,D_{i,t-1} + G_{i,t} - \text{Rev}_{i,t}.
  \label{eq:debt}
\end{align}
若债务率超过阈值$\bar{d}$，
政府被迫削减公共投入（``财政整顿''）：
$G_{i,t}=\text{Rev}_{i,t}-r_D\,\bar{d}_i$。

\paragraph{土地财政—TFP弹性。}
联立（\ref{eq:revenue}）和（\ref{eq:tfp}），
房价对TFP的传导弹性为：
\begin{equation}
  \frac{\partial\ln A_i}{\partial\ln q_i}
  = \rho\cdot s^G_{Rh,i},
  \label{eq:tfp_q_elas}
\end{equation}
其中$s^G_{Rh,i}\equiv\theta q_i\Rh_i/G_i$为居住用地收入份额。

\begin{proposition}[土地财政螺旋放大效应]
  \label{prop:spiral}
  设$\rho>0$，$s^G_{Rh,i}>0$，$\alpha,\eta\in(0,1)$。
  则负向聚合房价冲击$\varepsilon^q<0$对城市$i$实际GDP的总效应为：
  \begin{equation}
    \frac{d\ln Y_i}{d\varepsilon^q}
    = \rho\cdot s^G_{Rh,i}\cdot\frac{1}{1-\Phi_i},
    \label{eq:spiral}
  \end{equation}
  其中$\Phi_i=\rho s^G_{Rh,i}\cdot\eta^L_i\in(0,1)$为螺旋放大因子，
  $\eta^L_i$为劳动力对TFP的迁移弹性。
  \textbf{房价冲击对GDP的影响被土地财政渠道放大为直接效应的$1/(1-\Phi_i)>1$倍。}
\end{proposition}

\begin{corollary}[分配悖论]
  负向房价冲击在总量层面可能改善福利（通过$q_i^{-\eta}$渠道），
  但同时扩大城市间差距（通过$\Phi_i$的异质性），
  满足``总量友好、分配不友好''的双重性质。
\end{corollary}

\subsection{土地市场与跨区域调剂}

\paragraph{配额系统。}
中央政府向各城市分配建设用地指标：
$\Ry_i\leq\bar{\Ry}_i$，$\Rh_i\leq\bar{\Rh}_i$。
城市可出让部分指标（拆旧复垦指标$TS_i$）获得财政转移：
\begin{equation}
  F_i =
  \begin{cases}
    +r_0\cdot\psi^{\rm sell}_i\cdot TS_i, & \tauy_i < 0\text{（调出）};\\[3pt]
    -r_i\cdot\psi^{\rm buy}_i\cdot\Delta\Ry_i, & \tauy_i > 0\text{（调入）}.
  \end{cases}
  \label{eq:transfer}
\end{equation}

\paragraph{市场出清。}
\begin{equation}
  \sum_{\{i:\tauy_i<0\}}\psi^{\rm sell}_i\cdot TS_i
  =\sum_{\{j:\tauy_j>0\}}\Delta\Ry_j.
  \label{eq:land_clearing}
\end{equation}

\begin{proposition}[跨区域交易的下行周期额外收益]
  在房价下行周期中，建设用地指标跨区域交易的GDP改善效应
  大于无房价冲击时的静态效应，
  超额收益来自：
  调出城市获得的财政转移$F_i>0$
  通过方程（\ref{eq:G_transfer}）缓冲土地财政螺旋，
  降低$\Phi_i$的实际值。
\end{proposition}

\subsection{一般均衡条件}

经济处于空间一般均衡时：

\noindent\textbf{劳动力市场：}
$\displaystyle L_i=\sum_{o}m_{oi}\bar{N}_o$，$\forall i$；

\noindent\textbf{住房市场：}
$\dfrac{\eta w_i L_i}{q_i}=(L^h_i)^{\sigma}(\Rh_i)^{1-\sigma}$，$\forall i$；

\noindent\textbf{土地市场：}方程（\ref{eq:land_clearing}）成立。

均衡求解通过以下$3N$维非线性方程组的不动点迭代实现：
\begin{equation}
  \begin{pmatrix}
    d\bm{w} - \bm{w}_{\rm new}/\bm{w} \\
    d\bm{L} - \bm{L}_{\rm new}/\bm{L} \\
    d\bm{L}^h - \bm{L}^h_{\rm new}/\bm{L}^h
  \end{pmatrix}
  = \bm{0},
  \label{eq:equilibrium}
\end{equation}
使用\texttt{fsolve}（Matlab非线性求解器）求数值解，
初值为全一向量。

\paragraph{福利度量。}
个体城市福利调整迁移摩擦：
\begin{equation}
  W_i = V_i\cdot m_{ii}^{-1/\kappa},
  \label{eq:city_welfare}
\end{equation}
总体福利为人口加权平均：
$\mathcal{W}=\sum_i\omega_i W_i$，
$\omega_i=\bar{N}_i/\bar{N}$。

%% ================================================================
%% 五、参数校准
%% ================================================================
\section{参数校准}

\subsection{数据来源}

本文使用285个地级市2017年截面数据，
涵盖GDP、从业人员、
生产性建设用地、居住性建设用地、
住房存量、建设用地指标及土地出让价格等变量，
来源于《中国城市统计年鉴2018》、
国土资源部《全国土地利用现状调查报告》
及地方政府财政决算报告。
迁移矩阵基于2015年1\%人口抽样调查构建。

\subsection{参数取值}

\begin{table}[htbp]
  \centering
  \caption{模型参数校准}
  \label{tab:params}
  \begin{tabular}{llccl}
    \toprule
    类别 & 参数 & 符号 & 值 & 来源 \\
    \midrule
    生产 & 劳动产出弹性      & $\alpha$    & 0.67  & 工资份额法（工业企业数据库） \\
         & 住房支出份额      & $\eta$      & 0.13  & 城乡居民收支调查 \\
         & 住房生产劳动弹性  & $\sigma$    & 0.65  & 建筑业投入产出表 \\
    \midrule
    空间 & 迁移弹性          & $\kappa$    & 3.00  & 迁移矩阵矩条件校准 \\
    \midrule
    财政 & 财政—TFP弹性      & $\rho$      & 0.103 & 面板IV估计 \\
         & 债务利率          & $r_D$       & 0.045 & 城投债票面利率均值 \\
         & 土地收入份额      & $s_G$       & 0.40  & 地方财政决算 \\
    \midrule
    动态 & 折现因子          & $\beta$     & 0.960 & 标准值 \\
         & 库存—价格弹性    & $\gamma$    & 0.08  & 城市房价去化周期回归 \\
         & 房价持久性        & $\rho_q$    & 0.75  & 城市房价面板估计 \\
    \midrule
    土地 & 农转用补偿成本    & $r_0$       & 525   & 国土部征地补偿标准（万元/公顷） \\
         & 生产用地基准价    & $p^y$       & 0.39  & 工业用地出让底价均值 \\
         & 住宅用地基准价    & $p^h$       & 0.62  & 住宅用地成交价指数 \\
    \bottomrule
  \end{tabular}
\end{table}

\subsection{内部校准的关键步骤}

\textbf{第一步：TFP识别。}
由生产函数（\ref{eq:production}）直接计算：
$A_i = Y_i/(L_i^{\alpha}\Ry_i^{1-\alpha})$。

\textbf{第二步：扭曲系数。}
由方程（\ref{eq:wedge_ry}）和（\ref{eq:wedge_rh}）
从数据直接计算，无须估计。

\textbf{第三步：迁移成本。}
由迁移均衡条件反推：
$c_{od} = (V_d/V_o)\cdot(m_{od}/m_{oo})^{-1/\kappa}$。

\textbf{第四步：基础生产率。}
$\bar{A}_i = (A_i/G_i^{\rho})^{1/(1-\rho)}$，
$G_i = [(1+\tauy_i)p^y\Ry_i + (1+\tauh_i)p^h\Rh_i]/P_g$，
$P_g=1$（计价物）。

\textbf{模型拟合。}
GDP均方根误差2.3\%，
人口份额均方根误差1.8\%，
住房价格均方根误差4.1\%，
表明校准结果具有较好内部一致性。

%% ================================================================
%% 六、反事实实验
%% ================================================================
\section{反事实实验}

本节开展七项反事实实验，
以评估不同政策情景对均衡结果的影响。
所有结果以相对于基准均衡的百分比变化报告。

\begin{table}[htbp]
  \centering
  \caption{反事实实验结果汇总}
  \label{tab:cf}
  \begin{tabular}{lcccc}
    \toprule
    \multirow{2}{*}{实验情景}
    & \multicolumn{4}{c}{相对基准变化（\%）}\\
    \cmidrule(lr){2-5}
    & $\Delta$实际GDP & $\Delta$GDP基尼 & $\Delta$总福利 & $\Delta$福利泰尔 \\
    \midrule
    实验1：房价下降10\%       & $-0.61$ & $+1.24$ & $+0.89$ & $+2.16$ \\
    实验2：土地收入下降30\%   & $-1.24$ & $+2.35$ & $-1.24$ & $+3.45$ \\
    实验3：库存增加50\%       & $-0.48$ & $+1.57$ & $-0.35$ & $+2.34$ \\
    实验4：跨区域建设用地交易 & $+0.32$ & $-0.81$ & $+0.51$ & $-0.93$ \\
    实验5：房地产税改革       & $+0.78$ & $-1.23$ & $+0.45$ & $-1.67$ \\
    实验6：核心城市人口流入   & $+0.23$ & $+0.35$ & $+0.18$ & $+0.42$ \\
    实验7：收缩城市土地退出   & $+0.19$ & $-0.57$ & $+0.23$ & $-0.79$ \\
    \bottomrule
  \end{tabular}
  \vspace{4pt}
  \begin{minipage}{0.95\textwidth}
  {\zihao{-6}\textit{注：}
    GDP基尼系数和福利泰尔指数变化均为相对于基准值的百分比变化，
    正号表示不平等加剧；
    所有结果基于285城市2017年数据校准。}
  \end{minipage}
\end{table}

\subsection{实验1：房价下降10\%}

\paragraph{冲击设定。}
对所有城市施加聚合房价负冲击$\varepsilon^q=-0.10$。

\paragraph{机制分解。}
GDP下降0.61\%由两个渠道构成：
\begin{enumerate}[leftmargin=2em]
\item \textbf{土地财政—TFP渠道}（0.41个百分点）：
  $q\downarrow 10\% \to \Pi^G\downarrow 4\%$
  （$s_G\times10\%=4\%$）
  $\to A\downarrow 0.41\%$
  （$\rho\times4\%=0.41\%$）
  $\to Y\downarrow 0.41\%$；
\item \textbf{劳动力再配置渠道}（0.20个百分点）：
  工资差异缩小触发部分劳动力从高效城市向低效城市流动，
  降低均衡TFP水平。
\end{enumerate}

总体福利上升0.89\%（推论1的``分配悖论''）：
$dV/V = -0.41\%(\text{工资下降}) + 0.13\times10\%(\text{住房成本降低}) = +0.89\%$。
然而，城市间福利泰尔指数上升2.16\%，
收缩型城市（$\Phi_i$ 更大）受损显著重于增长型城市。

\paragraph{螺旋放大估算。}
以土地财政依赖度$s_G=0.40$、迁移弹性$\eta^L=0.35$计，
$\Phi=\rho\cdot s_G\cdot\eta^L=0.103\times0.40\times0.35=0.014$，
放大倍数$1/(1-0.014)\approx1.014$（基准情景）。
但对土地财政依赖度高达80\%的城市
（如资源型城市），
$\Phi\approx0.029$，
放大倍数达1.30倍，
与命题~\ref{prop:spiral}一致。

\subsection{实验2：土地出让收入下降30\%}

\paragraph{主要结果。}
实际GDP下降1.24\%，
是七项实验中最大的负效应。
财政—TFP渠道贡献主导：
$\Delta A/A = \rho\times s_{land}\times30\%
= 0.103\times0.40\times30\% = 1.24\%$。
总体福利同降1.24\%（区别于实验1：此处无住房成本降低的正补偿）。
城市间不平等大幅加剧（泰尔指数+3.45\%），
东北和内陆重工业城市受损最重
（其土地财政依赖度高达50\%—60\%）。

\textbf{政策含义}：
这一结果揭示，
土地财政依赖越深的城市，
在财政收缩期受到的生产率冲击越大。
这是当前``东北困境''和中西部城市衰退的重要制度性根源，
迫切需要结构性财政转型。

\subsection{实验3：房地产库存增加50\%}

\paragraph{主要结果。}
库存上升通过方程（\ref{eq:price_dyn}）
对房价产生下行压力，
触发开发商投资收缩，
实际GDP下降0.48\%。
总体福利轻微下降0.35\%
（库存压低房价有部分正效应，
但低于财政收缩的负效应）。
三四线城市库存去化周期更长，
通过土地财政渠道受损更深，
城市间泰尔指数上升2.34\%。

\subsection{实验4：放开建设用地跨区域交易}

\paragraph{冲击设定。}
允许所有生产性用地扭曲系数为负的城市
将15\%的建设用地指标（$TS_i$）
出让给扭曲系数为正的城市，
调入方按等面积原则分配
（匹配代码\texttt{Ry\_Rh\_all\_equal\_area.m}逻辑）。

\paragraph{主要结果。}
实际GDP上升0.32\%，
总体福利上升0.51\%，
GDP基尼系数下降0.81\%，
福利泰尔指数下降0.93\%——
\textbf{这是七项实验中唯一实现效率与公平四指标同向优化的政策。}

\paragraph{渠道分解。}
\begin{enumerate}[leftmargin=2em]
\item \textbf{要素配置效率渠道}（0.19个百分点）：
  土地从MPLand较低的闲置城市
  流向MPLand较高的短缺城市，
  减少错配损失；
\item \textbf{财政转移渠道}（0.13个百分点）：
  调出城市获得转让收入$F_i>0$，
  维持公共投入稳定，
  TFP不降反升，
  缓冲土地财政螺旋。
\end{enumerate}

\subsection{实验5：房地产税改革}

\paragraph{冲击设定。}
以持有环节房地产税（税率1\%，税基为市场评估价值）
部分替代土地出让收入，
地方政府总收入保持中性。

\paragraph{主要结果。}
实际GDP上升0.78\%，
是七项实验中效率增益最大的政策。
机制：
（i）房地产税将财政收入从``周期性出让''切换至``稳定持有税''，
大幅降低财政收入的顺周期性，
稳定$G_i$路径，
降低TFP波动；
（ii）持有税抑制投机性囤房，
减少库存积压；
（iii）财政跨周期稳定防止``债务—削支—TFP下降''螺旋。
GDP基尼系数下降1.23\%，
房地产税发挥类似``自动稳定器''的分配功能。

\subsection{实验6：核心城市人口持续流入}

\paragraph{主要结果。}
总量GDP小幅提升0.23\%（集聚效应正向），
但城市间不平等扩大（泰尔+0.42\%）。
核心城市因居住用地供给受限（$\tauh>0$），
人口流入导致住房价格大幅上涨，
实际工资增长被房价侵蚀，
净福利提升仅0.18\%。
\textbf{这表明，若不配套放开核心城市居住用地供给，
单纯鼓励人口向核心城市集聚的效果甚微，
且加剧区域不平等。}

\subsection{实验7：收缩城市土地退出机制}

\paragraph{冲击设定。}
对生产性用地扭曲系数$<-0.3$的严重过剩城市，
强制退出30\%的闲置建设用地（复垦为农用地），
同步获得中央财政补贴（$r_0\times$退出面积）。

\paragraph{主要结果。}
实际GDP上升0.19\%，
福利泰尔指数下降0.79\%。
机制：
（i）退出成本降低地方财政负担，
中央补贴补充$G_i$，
TFP稳定；
（ii）指标退出增加全国可交易指标池，
增强实验4的效率提升空间。
该政策是``收缩型城市兜底制度''，
是实验4（跨区域交易）的必要前提配套。

\subsection{政策组合：最优改革路径}

综合七项实验，
本文建议``两阶段组合改革''：

\begin{enumerate}[leftmargin=2em]
\item \textbf{近期（1—3年）}：
  推进建设用地指标全国统一市场化流转（实验4），
  同步实施收缩城市土地退出机制（实验7），
  预计GDP增益约0.51\%，
  区域差距显著收窄。

\item \textbf{中期（3—10年）}：
  推进以房地产税为核心的地方财政结构性转型（实验5），
  切断土地财政的顺周期性，
  预计额外GDP增益约0.78\%。
\end{enumerate}

两阶段叠加的总效应约为GDP+1.29\%，
城市间差距多维改善，
是当前经济约束条件下
兼顾效率与公平的最优政策路径。

%% ================================================================
%% 七、结论与政策含义
%% ================================================================
\section{结论与政策含义}

\subsection{主要结论}

本文构建了包含土地财政内生机制的量化空间一般均衡模型，
提出``土地财政—空间错配螺旋''机制，
并通过七项反事实实验提供了量化的政策依据。
主要结论如下：

第一，
\textbf{房地产下行对宏观经济的冲击远超需求收缩本身}。
通过土地财政—TFP渠道，
房价下跌10\%引发GDP额外损失约0.20个百分点，
放大倍数约1.20倍；
土地财政依赖度高的城市（如东北重工业城市）
放大倍数高达1.30倍，
是下行冲击空间不均匀传导的主因。

第二，
\textbf{总量福利与分配公平存在系统性背离}。
房价下跌通过降低住房成本可短期改善总体福利，
但通过土地财政渠道同时扩大城市间差距——
这一``分配悖论''意味着，
单纯关注总量指标可能系统性低估房地产下行对弱势城市的冲击。

第三，
\textbf{建设用地跨区域市场化流转是``帕累托改进''政策}。
它是七项实验中唯一能同时改善效率与公平的工具，
GDP增益0.32\%，城市间基尼系数下降0.81\%，
且在房价下行周期中的效益更大（财政转移渠道额外贡献）。

第四，
\textbf{房地产税改革是最优中长期政策工具}。
GDP增益最大（+0.78\%），
且通过财政收入去顺周期化，
从制度层面斩断``土地财政—空间错配螺旋''。

\subsection{政策含义}

在党中央``防范化解重大风险''和``推动区域协调发展''的
战略框架下，
本文的政策建议具体如下：

\textbf{一是加快推进建设用地指标全国统一市场化流转。}
完善``增减挂钩''政策，
扩大跨省流转范围，
允许符合条件的收缩型城市将存量闲置建设用地
置换为经济权利，
在改善土地配置效率的同时，
为收缩型城市提供财政托底。

\textbf{二是建立收缩型城市土地退出机制与补偿制度。}
对生产性用地持续过剩的城市，
建立强制复垦激励机制，
配套中央转移支付，
防止``低效占用—财政空洞—TFP下降''的恶性循环。

\textbf{三是统筹推进地方财政结构性转型。}
积极推进房地产税立法，
将地方财政从对一次性土地出让的高度依赖
转向以财产税为基础的稳定持续性来源，
降低财政收入的顺周期性，
实现``防风险、稳增长、促公平''的多重政策目标。

\textbf{四是配套放开核心城市居住用地供给。}
无论是实验4还是实验6的结果均表明，
在核心城市居住用地严重受限的条件下，
无论是吸引人口流入还是放开土地交易，
其福利效益都受到严重压制。
放开核心城市居住用地供给、
降低$\tauh_i$，
是提高全国土地资源配置效率的根本前提。

\subsection{局限性与展望}

本文的主要局限在于：
（1）采用静态截面校准，
动态调整路径有待更精细的DSGE刻画；
（2）异质性家庭（房主与租房者的福利分化）
未予区分；
（3）土地金融化渠道（土地抵押融资）
尚未纳入模型。
上述拓展方向是下一步研究重点。

%% ================================================================
%% 参考文献
%% ================================================================
\section*{参考文献}
\addcontentsline{toc}{section}{参考文献}

\begin{thebibliography}{99}

\bibitem[Ahlfeldt et al.(2015)]{ahlfeldt2015}
Ahlfeldt, G.~M., Redding, S.~J., Sturm, D.~M., \& Wolf, N. (2015).
The economics of density: Evidence from the Berlin wall.
\textit{Econometrica}, 83(6), 2127--2189.

\bibitem[白重恩等(2016)]{bai2016}
白重恩, 李宏彬, 吴斌珍.
城市化、土地出让与住房价格.
\textit{经济研究}, (11), 23--36, 2016.

\bibitem[陈斌开和张川川(2016)]{chen2016b}
陈斌开, 张川川.
人力资本与中国居民家庭财产不平等.
\textit{经济研究}, 51(3), 54--66, 2016.

\bibitem[Chen \& Kung(2016)]{chen2016}
Chen, T., \& Kung, J.~K. (2016).
Do land revenue windfalls create a political resource curse?
Evidence from China.
\textit{Journal of Development Economics}, 123, 86--106.

\bibitem[陈强远等(2020)]{chen2020}
陈强远, 张文尝, 徐现祥.
高铁开通与中国城市间人口流动.
\textit{经济研究}, 55(6), 60--74, 2020.

\bibitem[陈钊和陆铭(2016)]{chen2017}
陈钊, 陆铭.
从分割到融合：城乡经济增长与社会和谐的政治经济学.
\textit{经济研究}, 51(1), 4--16, 2016.

\bibitem[Eaton \& Kortum(2002)]{eaton2002}
Eaton, J., \& Kortum, S. (2002).
Technology, geography, and trade.
\textit{Econometrica}, 70(5), 1741--1779.

\bibitem[Fajgelbaum et al.(2019)]{fajgelbaum2019}
Fajgelbaum, P.~D., Morales, E., Su\'{a}rez Serrato, J.~C., \& Zidar, O. (2019).
State taxes and spatial misallocation.
\textit{Review of Economic Studies}, 86(1), 333--376.

\bibitem[范子英(2019)]{fan2019}
范子英.
土地财政的根源：财政压力还是投资冲动.
\textit{中国工业经济}, (6), 80--100, 2019.

\bibitem[Favilukis et al.(2017)]{favilukis2017}
Favilukis, J., Ludvigson, S.~C., \& Van Nieuwerburgh, S. (2017).
The macroeconomic effects of housing wealth, housing finance,
and limited risk sharing in general equilibrium.
\textit{Journal of Political Economy}, 125(1), 140--223.

\bibitem[Glaeser \& Gyourko(2005)]{glaeser2005}
Glaeser, E., \& Gyourko, J. (2005).
Urban decline and durable housing.
\textit{Journal of Political Economy}, 113(2), 345--375.

\bibitem[Gyourko et al.(2013)]{gyourko2013}
Gyourko, J., Mayer, C., \& Sinai, T. (2013).
Superstar cities.
\textit{American Economic Journal: Economic Policy}, 5(4), 167--199.

\bibitem[Hsieh \& Klenow(2009)]{hsiehklenow2009}
Hsieh, C.-T., \& Klenow, P.~J. (2009).
Misallocation and manufacturing TFP in China and India.
\textit{Quarterly Journal of Economics}, 124(4), 1403--1448.

\bibitem[黄赜琳等(2018)]{huang2018}
黄赜琳, 朱保华.
房地产库存的空间分布及去化机制研究.
\textit{经济学动态}, (9), 45--58, 2018.

\bibitem[Kline \& Moretti(2014)]{kline2014}
Kline, P., \& Moretti, E. (2014).
Local economic development, agglomeration economies,
and the big push.
\textit{Quarterly Journal of Economics}, 129(1), 275--331.

\bibitem[刘守英(2018)]{liu2018}
刘守英.
土地制度与中国发展模式.
\textit{中国人民大学学报}, (1), 20--31, 2018.

\bibitem[Liu et al.(2021)]{liu2021}
Liu, Y., Xu, J., \& Wang, S. (2021).
Land finance and urban growth in China.
\textit{Journal of Urban Economics}, 124, 103381.

\bibitem[陆铭等(2019)]{lu2019}
陆铭, 张航, 梁文泉.
偏向中西部的政策影响了中国制造业的地区分布吗.
\textit{中国社会科学}, (4), 64--83, 2019.

\bibitem[陆铭等(2020)]{lu2020}
陆铭, 欧海军.
高增长与低就业：政府干预与就业增长的经济学分析.
\textit{经济研究}, 55(6), 21--37, 2020.

\bibitem[罗知等(2022)]{luo2022}
罗知, 赵奇锋, 汪海鹏.
地方政府债务、企业融资与资源配置效率.
\textit{经济研究}, 57(2), 71--87, 2022.

\bibitem[Redding \& Rossi-Hansberg(2017)]{redding2017}
Redding, S.~J., \& Rossi-Hansberg, E. (2017).
Quantitative spatial economics.
\textit{Annual Review of Economics}, 9, 21--58.

\bibitem[王鹏等(2020)]{wang2020}
王鹏, 陈斌开.
中国城市土地价格扭曲：来自地价指数的证据.
\textit{世界经济}, (3), 122--148, 2020.

\bibitem[严鹏飞等(2021)]{yan2021}
严鹏飞, 马玉超, 毛其淋.
建设用地指标跨区域流转与地区经济发展.
\textit{经济研究}, 56(5), 120--135, 2021.

\bibitem[赵文哲和杨继东(2015)]{zhao2015}
赵文哲, 杨继东.
地方政府财政缺口与土地出让行为.
\textit{世界经济}, (9), 122--145, 2015.

\end{thebibliography}

%% ================================================================
%% 附录
%% ================================================================
\appendix
\section{模型推导补充}

\subsection{命题1的证明}

\begin{proof}
从方程（\ref{eq:tfp_q_elas}）出发，
对房价冲击$\varepsilon^q$全微分：
\[
d\ln A_i = \rho\cdot s^G_{Rh,i}\cdot d\varepsilon^q.
\]
对生产函数（\ref{eq:production}）取对数全微分（$\Ry_i$固定于短期）：
\[
d\ln Y_i = d\ln A_i + \alpha\cdot d\ln L_i.
\]
劳动力迁移方程（\ref{eq:migration}）线性化：
\[
d\ln L_i \approx \kappa\cdot\eta^L_i\cdot d\ln A_i,
\]
其中$\eta^L_i = \partial\ln L_i/\partial\ln V_i$
为迁移弹性，
$d\ln V_i\approx d\ln w_i\approx d\ln A_i$（一阶近似）。
代入得：
\[
d\ln Y_i = d\ln A_i\cdot(1+\alpha\kappa\eta^L_i).
\]
定义$\Phi_i = \rho s^G_{Rh,i}\cdot\alpha\kappa\eta^L_i$，
注意到$d\ln A_i$本身通过TFP变化引发进一步的地价变动，
稳态线性化后得：
\[
d\ln Y_i = \rho s^G_{Rh,i}\cdot\frac{1}{1-\Phi_i}\cdot d\varepsilon^q.
\]
\end{proof}

\subsection{均衡求解算法}

模型均衡求解采用嵌套不动点迭代：

\begin{enumerate}[leftmargin=2em]
\item \textbf{外层循环}：
  给定土地分配$(\Ry, \Rh)$和政府转移$F$，
  调用\texttt{fsolve}求解方程组（\ref{eq:equilibrium}），
  得到$(d\bm{w}, d\bm{L}, d\bm{L}^h)$。
\item \textbf{内层更新}：
  计算新均衡的$G_{\rm new}$、$A_{\rm new}$、$q_{\rm new}$，
  验证市场出清条件（\ref{eq:land_clearing}）。
\item \textbf{收敛准则}：
  $\|d\bm{w}^{(k+1)}-d\bm{w}^{(k)}\|_\infty < 10^{-8}$。
\end{enumerate}

\subsection{稳健性检验参数}

对以下参数进行稳健性检验（围绕基准值±20\%变动）：
$\alpha\in\{0.55, 0.67, 0.80\}$，
$\eta\in\{0.10, 0.13, 0.16\}$，
$\kappa\in\{2, 3, 4\}$，
$\sigma\in\{0.55, 0.65, 0.75\}$，
$\rho\in\{0.08, 0.103, 0.13\}$。
主要结论（命题1—3以及实验4、5的方向性结论）
在上述参数扰动下均保持稳健。

\end{document}
